% Introduction du chapitre II (13 août – 24 septembre 1916 ; ff. 10 r° – 25 r°), fragment LaTeX sans préambule.

À la mi-août 1916, Hippolyte Combes est à La Croix-du-Bac (Steenwerck), entre la Lys et le front d'Armentières, avec la 152nd Company de l'\textit{Army Service Corps}, qui ravitaille la 54th Brigade de la 18th (Eastern) Division. En six semaines, il passe de la Lys au Ternois, où la brigade s'exerce à l'attaque pendant quinze jours, puis gagne en trois étapes le plateau picard, à l'arrière de l'Ancre, où la division se prépare à attaquer Thiepval. Le service ne change pas ; l'interprète, lui, est désormais admis à la table des officiers. Le chapitre s'arrête la veille de son départ pour la première des deux permissions du carnet.

\subsection*{Le carnet}

Les douze derniers jours en Flandre sont faits d'attente (ff.~10~r\textsuperscript{o}–13~v\textsuperscript{o}). Hippolyte Combes va aux champs chercher le fourrage vert avec le lieutenant Curtis, qu'il juge durement : «~ce crétin de Curtis~» (f.~11~r\textsuperscript{o}). Il se promène souvent le long de la Lys avec Maurin, sans doute un autre interprète (ff.~10~r\textsuperscript{o}, 11~r\textsuperscript{o}, 12~v\textsuperscript{o}, 13~v\textsuperscript{o}). Le 19, il accompagne Curtis à cheval dans la «~zone exposée, derrière ligne de feu~» : Fleurbaix, puis Armentières, aux façades trouées, presque vidée de ses habitants (ff.~11~r\textsuperscript{o}–v\textsuperscript{o}). Le 20, il prend café et rhum chez M.~Barbry avec deux adjudants interprètes, Lupin et Schumacher (f.~12~r\textsuperscript{o}). Le 24, il trouve dans l'église de Steenwerck une «~relique d St~Hippte, mon patron !~» et note le soir son «~Impression de melancolie au moment de m'eloigner de ce village pour toujours~» (f.~13~v\textsuperscript{o}).

Le 25~août, la compagnie embarque ses chevaux et ses voitures à Bailleul et gagne Saint-Pol par Hazebrouck, Berguette et Lillers (f.~14~r\textsuperscript{o}). Dans le wagon, il note des «~Reflexions amères~» sur l'État qui, après les corps, prend aussi l'intelligence et le cœur (f.~14~v\textsuperscript{o} ; lecture en partie incertaine). À Marquay, «~petit hameau perdu en haut d'1 plateau d Artois~», il trouve à grand-peine une «~chambrette exiguë~» dans une maison «~très mal tenue, sale, moisie~» (ff.~14~r\textsuperscript{o}–15~r\textsuperscript{o}). Il pleut presque chaque jour, et le fourrage, matin et soir, dans les champs détrempés, fait une «~Vie vegetative deprimante \& abrutissante~» (ff.~15~v\textsuperscript{o}–16~r\textsuperscript{o}). Le 29, il note l'entrée en guerre de la Roumanie (f.~15~v\textsuperscript{o}). Le 4~septembre, ses paquets faits, il se réjouit de quitter «~cet obscur \& boueux hameau~» ; le lendemain, le départ est retardé : «~encore 8 jours à moisir ici, hélas!~» (f.~18~v\textsuperscript{o}). Avec Curtis, il rend visite à une fermière, Mme~Vauchel (f.~19~r\textsuperscript{o}).

Le 9, la compagnie part après le défilé de «~3 brigades division~» et fait étape à Monts-en-Ternois ; le capitaine West invite Hippolyte Combes à «~messer~» avec les officiers (f.~20~r\textsuperscript{o}). Le 10, à Oppy, hameau de Beaudricourt, il va «~on duty, chz le maire~» (f.~20~v\textsuperscript{o}). Le 11, par Lucheux et Thièvres, la compagnie arrive à Raincheval, «~très gros village [...] fourmillant de tommies~» (f.~21~r\textsuperscript{o}). Le 12, il devient membre du mess de West «~à titre definitif~» : il y trouve la conversation anglaise et une «~sorte de club où me retirer pour lire revues \& journaux~» (f.~21~v\textsuperscript{o}). Le bureau de la Mission française, à Acheux, est une «~maisonnette toile \& planches~» (f.~22~r\textsuperscript{o}). L'attente continue : fourrage vert, promenades, visite à un collectionneur avec un autre interprète, culte wesleyen dans une petite salle en face de l'église (f.~22~v\textsuperscript{o} ; cf. \textit{Feuillets de Route}, chap.~\textsc{xix} : \textit{«~Sunday service in village church\LJinc{way}~»}, et le poème wesleyen copié par Hippolyte Combes, annexe, p.~\pageref{annexe:doc4} et~\pageref{annexe:doc6}). Le 18, appelé à la Mission, il la cherche à Hédauville, puis la trouve à Acheux, où le quartier général de la division est resté (f.~23~r\textsuperscript{o}) ; le 21, il y touche sa solde, 78~F~75 (f.~23~v\textsuperscript{o}).

Le 22, il apprend du capitaine «~Pant~» que «~we are moving demain~» : «~gde deception~». Sa permission est fixée au 25 (f.~24~r\textsuperscript{o}). Le 23, la compagnie quitte Raincheval ; près de Léalvillers, le lieutenant Pizzi arrive avec un message, et il faut «~tourner bride~» ; l'après-midi, la compagnie campe «~in the open~» près de Forceville, où Hippolyte Combes se fait une hutte avec une bâche verte (f.~24~v\textsuperscript{o}). Dernière ligne du chapitre : «~Canon, grondé tout le long du front~» (f.~25~r\textsuperscript{o}).

\subsection*{Le contexte}

Le secteur de la Lys passe pour calme depuis l'attaque de Fromelles (voir le chapitre~\textsc{i}), mais d'un calme relatif : dans la nuit du 14 au 15~août, un raid simulé du 7th Bedfordshire devant Armentières tue un sous-lieutenant de dix-huit ans et trois hommes\footnote{Bedfordshire Archives Service, X550/8/1/1608, \textit{War Diary}, 7th (Service) Battalion Bedfordshire Regiment, août 1916, 15~août (ci-après JMO 7/Bedf) ; CWGC, Second Lieutenant David Arnold Roeber, Ration Farm Military Cemetery, La Chapelle-d'Armentières.}. Le carnet ne dit rien de la visite de George~V, du prince de Galles et du général Plumer, le 14~août, à l'école de grenadiers des Australiens, à Sailly-sur-la-Lys, à deux kilomètres et demi de La Croix-du-Bac\footnote{A.~D.~Ellis, \textit{The Story of the Fifth Australian Division}, Londres, Hodder \& Stoughton, [1920], p.~120–121.}. Les «~Ecossais~» ajoutés après coup à la description d'Armentières (f.~11~v\textsuperscript{o}) sont les Highlanders de la 152nd Brigade (51st Division), arrivés les 15 et 16~août dans une ville encore habitée par de nombreux civils\footnote{F.~W.~Bewsher, \textit{The History of the 51st (Highland) Division, 1914–1918}, Édimbourg et Londres, Blackwood, 1921, p.~87–88.}. Vers le 22~août, la 34th Division, revenue de la Somme, relève la 18th\footnote{J.~Shakespear, \textit{The Thirty-Fourth Division, 1915–1919}, Londres, Witherby, 1921, p.~74–75.}. L'infanterie de la 54th Brigade embarque à la gare ouest de Bailleul dans la nuit du 24 au 25 ; chevaux et voitures suivent le lendemain. Elle cantonne ensuite à La Thieuloye, Marquay, Orlencourt et Bailleul-aux-Cornailles\footnote{JMO 7/Bedf, ordre d'opérations n\textsuperscript{o}~26 du 23~août 1916.}.

Le Ternois est l'arrière de la Third Army, qui tient le front d'Hébuterne à Vimy. La 18th Division y reçoit «~three weeks' battle training~» avant de retourner sur la Somme\footnote{W.~Miles (comp.), \textit{Military Operations, France and Belgium, 1916}, vol.~II, Londres, Macmillan, 1938 (histoire officielle britannique, ci-après \textit{OH 1916}, II), p.~394 et p.~543, n.~3 ; G.~H.~F.~Nichols, \textit{The 18th Division in the Great War}, Édimbourg et Londres, Blackwood, 1922, p.~78.}. Les «~manœuvres speciales~» vues près de Ligny (f.~17~r\textsuperscript{o}) en font partie : attaque au bois des Hérombus le 28~août, attaque à l'aube contre deux lignes de tranchées le 31, devant les commandants du corps et de la division, puis creusement et progression derrière un barrage\footnote{JMO 7/Bedf, 28 et 31~août 1916 ; \textit{ibid.}, X550/8/1/1609, septembre 1916, 5–8~septembre.}. À Chelers, le barrage roulant est figuré par une ligne d'hommes portant des fanions bleus\footnote{T.~M.~Banks et R.~A.~Chell, \textit{With the 10th Essex in France}, 2\textsuperscript{e}~éd., Londres, Gay \& Hancock, 1924, p.~144–145.}. L'école de mortiers de tranchée de Ligny est celle de la Third Army, que le roi avait visitée le 9~août\footnote{Imperial War Museum, photographie Q~961 (Lt~E.~Brooks), 9~août 1916.}. Le journal du 7th Bedfordshire, à La Thieuloye, donne presque jour par jour le même temps que le carnet, ce qui confirme ses dates : «~Very heavy rain~» le 26~août, «~Heavy thunderstorm~» le 29, «~Wet day~» le 30, «~Fine day~» le 31\footnote{JMO 7/Bedf, 26–31~août 1916.}. La Roumanie déclare la guerre à l'Autriche-Hongrie le 27~août, l'Allemagne à la Roumanie le 28 : d'où les «~Confirmations~» du 29\footnote{C.~Falls (comp.), \textit{Military Operations, Macedonia}, vol.~I, Londres, HMSO, 1933, p.~139.}. Le départ retardé du 5~septembre correspond à l'annulation, ce matin-là à 5~heures, de l'ordre qui faisait passer la division à la Reserve Army ; le transfert eut lieu dans la nuit du 10 au 11\footnote{JMO 7/Bedf, ordres n\textsuperscript{os}~28 (4~septembre, annoté «~Cancelled 5 a.m. 5.9.16~») et 29 (8~septembre).}.

Passée au II Corps de la Reserve Army du général Gough, la division envoie de Raincheval ses officiers reconnaître le terrain devant Thiepval\footnote{JMO 7/Bedf, 12–13~septembre 1916 ; Nichols, \textit{op.~cit.}, p.~78.}. Le 15~septembre, plus au sud, s'ouvre la bataille de Flers-Courcelette, première action des chars ; le même jour, la 54th Brigade doit se tenir «~ready to move at 6 hours notice~»\footnote{\textit{OH 1916}, II, p.~xvii ; JMO 7/Bedf, 15~septembre 1916.}. Les villages sont surpeuplés : selon l'histoire officielle, cet arrière-front, où il fallait loger plus de 400\,000~hommes et 100\,000~chevaux, n'avait que des routes rares et médiocres\footnote{J.~E.~Edmonds (comp.), \textit{Military Operations, France and Belgium, 1916}, vol.~I, Londres, Macmillan, 1932, p.~271.}. Le quartier général de la division reste à Acheux jusqu'au 25~septembre, puis passe à Hédauville\footnote{Nichols, \textit{op.~cit.}, p.~79.}. Le 23, le 7th Bedfordshire va de Raincheval à Varennes par Léalvillers et Acheux, dans le mouvement qui conduit la brigade devant Thiepval\footnote{JMO 7/Bedf, 23~septembre 1916.}. La préparation d'artillerie contre la crête de Thiepval commence ce jour-là ; désormais «~a thick mist always prevailed in the mornings and evenings~» : c'est l'«~épais brouillard~» du 24 (f.~25~r\textsuperscript{o}), et le canon de ce soir-là est celui de cette préparation\footnote{\textit{OH 1916}, II, p.~372 et 393–394.}.

Le cantonnement (voir le chapitre~\textsc{i}) passe par le maire, intermédiaire légal du logement des troupes : d'où le service «~chz le maire~» de Beaudricourt\footnote{Fernand-Laurent, \textit{Chez nos alliés britanniques. Notes et souvenirs d'un interprète}, Paris, Boivin, 1917, p.~225–234 ; loi du 3~juillet 1877 relative aux réquisitions militaires.}. Le fourrage vert (voir le chapitre~\textsc{i}) se paie aux cultivateurs\footnote{«~Claims~», par «~A Divisional Claims Officer~», \textit{Blackwood's Magazine}, vol.~201, n\textsuperscript{o}~1220, juin 1917, p.~933.}. L'usage fait d'abord de l'interprète «~l'invité du mess~»\footnote{Fernand-Laurent, \textit{op.~cit.}, p.~64–65.}. Le capitaine West est Edward Maurice West, qui commande encore la 152nd Company en 1918\footnote{The National Archives (Kew), WO~339/21191 ; Nichols, \textit{op.~cit.}, p.~110 et 432.}. Schumacher est vraisemblablement le «~Schuhmacher, the interpreter at Brigade H.Q.~» de la 53rd Brigade, que l'historique du 10th Essex retrouve en 1918\footnote{Banks et Chell, \textit{op.~cit.}, p.~310.}. Curtis, Maurin et le capitaine «~Pant~» restent à identifier ; c'est sans doute le même «~Cap. Pant~» à qui Hippolyte Combes paiera sa part du mess de septembre, 29~F (f.~34~v\textsuperscript{o})\footnote{«~Pant~» abrège peut-être un nom ; aucun officier de l'\textit{Army Service Corps} nommé Panton ne figure dans les dossiers d'officiers des National Archives (Kew), séries WO~339 et WO~374 (recherche dans le catalogue \textit{Discovery}).}. Les prix montent avec la présence des troupes : le litre de bière, 30~c en Flandre (f.~7~v\textsuperscript{o}), coûte 50~c à Raincheval (f.~21~r\textsuperscript{o}). Si la lecture «~1~f~» est bonne, le quart de livre de chocolat Poulain revient à 8~F le kilo, quand le chocolat ordinaire valait 3,19~F à l'économat d'une compagnie de chemins de fer en juillet 1916\footnote{L.~March, \textit{Mouvement des prix et des salaires pendant la guerre}, Paris, PUF, et New Haven, Yale University Press, 1925, p.~199–200.}. \textit{Canada in Flanders} est le premier volume de l'histoire officielle du corps expéditionnaire canadien, vendu un shilling\footnote{Sir Max Aitken, \textit{Canada in Flanders. The Official Story of the Canadian Expeditionary Force}, t.~I, Londres, Hodder \& Stoughton, 1916 ; le t.~II porte une préface du 27~janvier 1917.}.

Le départ en permission, le 25~septembre, ouvre le chapitre~\textsc{iii}.
